% Authoritative LaTeX chapter. Edit this file for future revisions.
% Initially migrated 2026-09-11; no Markdown import occurs during PDF builds.
\chapter{상세 실행 예시와 LUT 제작}
\label{chap:04_examples}
이 부록의 기본 예시를 실행한 뒤 \hyperref[chap:09_cookbook]{부록~\ref*{chap:09_cookbook}의 목적별 예제 모음}에서 각도·대기·수질·자료 생성 조건을 확장한다.
목차 · \hyperref[chap:03_options]{옵션 사전} · \hyperref[chap:05_output_formats]{출력 해석}

\Needspace{17\baselineskip}
\section{공통 준비와 계산 모드}
\label{sec:auto:04_examples:1}
이 부록의 Bash 명령은 \textbf{\texttt{code/\allowbreak{}OCRT\_\allowbreak{}C}에서} 실행한다. \texttt{.\allowbreak{}/\allowbreak{}build/\allowbreak{}ocrt}는 현재 소스로 빌드한 파일이어야 한다. 먼저 출력 폴더와 기체흡수 off 인자를 준비한다.


\par\noindent\begin{minipage}{\linewidth}
\begin{ManualCode}
export OMP_NUM_THREADS=1
mkdir -p ../../task/manual_examples
GAS_OFF=(--gas-column-h2o 0 --gas-column-o3 0 --gas-column-no2 0
         --gas-column-o2 0 --gas-column-co2 0 --gas-column-ch4 0)
\end{ManualCode}
\end{minipage}\par

\texttt{GAS\_\allowbreak{}OFF}는 Bash 배열이다. 다른 셸에서 그대로 쓰지 않는다. PowerShell은 \hyperref[sec:example_first_run]{부록~\ref*{sec:example_first_run}의 \texttt{\$gasOff} 배열}을 만들고 \texttt{"\$\{GAS\_\allowbreak{}OFF[@]\}"}를 \texttt{@gasOff}, 실행파일을 \texttt{\& .\allowbreak{}\textbackslash{}build\textbackslash{}ocrt\_\allowbreak{}x86-\allowbreak{}64-\allowbreak{}v3.\allowbreak{}exe}, 줄연결 \texttt{\textbackslash{}}를 백틱으로 바꾼다. \hyperref[chap:examples_readme]{예시 도구}를 쓰면 두 운영체제에서 같은 사례 계획을 생성할 수 있다.


{\TableFont
\begin{longtable}{@{}>{\raggedright\arraybackslash}p{\dimexpr 0.3400\linewidth-4.00pt\relax}>{\raggedright\arraybackslash}p{\dimexpr 0.6600\linewidth-4.00pt\relax}@{}}
\toprule
\textbf{입력 구성} & \textbf{동작} \\
\midrule
\endfirsthead
\toprule
\textbf{입력 구성} & \textbf{동작} \\
\midrule
\endhead
\bottomrule
\endfoot
\texttt{-\allowbreak{}-\allowbreak{}vza V -\allowbreak{}-\allowbreak{}raa R} 모두 지정 & 한 방향 계산; 결과는 표준출력 \\
VZA·RAA 모두 생략하고 \texttt{-\allowbreak{}-\allowbreak{}output-\allowbreak{}full-\allowbreak{}grid F.\allowbreak{}csv} 지정 & 균일 VZA/RAA 격자 CSV \\
\texttt{-\allowbreak{}-\allowbreak{}batch-\allowbreak{}full-\allowbreak{}grid jobs.\allowbreak{}csv} & 해양 전체 격자 사례를 행 단위로 실행 \\
\end{longtable}
}

VZA·RAA를 하나만 지정하거나 단일 방향과 격자 옵션을 섞으면 오류다. 파일 경로는 명시한다. 특히 해양 격자는 \texttt{-\allowbreak{}-\allowbreak{}output-\allowbreak{}full-\allowbreak{}grid}를 명시해야 명시적 격자 경로의 수중 수치 기본값도 고정된다. \textbf{기본 2.5\textdegree{} 격자는 첫 동작 확인에 불필요하게 크므로 작은 격자부터 시작한다.}

\Needspace{17\baselineskip}
\section{순수 분자 산란: 검은 경계}
\label{sec:auto:04_examples:2}
해면 기여를 제외한 대기 기준값이다. 6기체 총량을 모두 0으로 지정하고 에어로졸을 생략한다.


\par\noindent\begin{minipage}{\linewidth}
\begin{ManualCode}
./build/ocrt --surface black --sza 40 --vza 30 --raa 90 \
  --wavelength 555 --pressure 1013.25 "${GAS_OFF[@]}" \
  > ../../task/manual_examples/ray_black.txt \
  2> ../../task/manual_examples/ray_black.log
\end{ManualCode}
\end{minipage}\par

\texttt{-\allowbreak{}-\allowbreak{}pressure 0}은 분자 산란을 없앤다. 기체 흡수를 끄는 옵션은 아니므로 이를 기체 off와 혼동하지 않는다. \texttt{-\allowbreak{}-\allowbreak{}surface black}의 LUT를 거친 바다 위 대기보정 LUT라고 부르지 않는다.

\Needspace{17\baselineskip}
\section{해양 대기보정용 Rayleigh LUT}
\label{sec:auto:04_examples:3}
\subsection{물리 조건부터 정하기}
\label{sec:auto:04_examples:3.1}
해양 대기보정에서 사용할 Rayleigh LUT의 한 구성은 \textbf{분자 산란 + 거친 Fresnel 해면 + 수중 광원 없음 + 에어로졸 없음 + 기체 흡수 없음}이다. 직달 태양광의 선글린트는 별도로 보정하는 경우가 많으므로 다음 예시는 \texttt{-\allowbreak{}-\allowbreak{}decouple-\allowbreak{}sunglint}로 해당 단일반사 항을 제외한다. 이것이 해면 경계의 모든 반사를 없애는 것은 아니다.

이 구성의 \texttt{rho\_\allowbreak{}I/\allowbreak{}Q/\allowbreak{}U}를 해당 정의의 Rayleigh 항으로 사용한다. 사용 중인 대기보정 시스템이 선글린트를 LUT에 포함한다면 \texttt{-\allowbreak{}-\allowbreak{}decouple-\allowbreak{}sunglint}를 제거하고, 그 정의를 LUT 메타데이터에 남긴다. 서로 다른 정의의 LUT를 섞으면 중복 보정 또는 누락이 발생한다.

\subsection*{글린트와 편광을 실제 계산 그림으로 확인하기}
그림~\ref{fig:measured-azimuth}는 동일한 태양·관측 천정각에서 바람과 직달 글린트 포함 여부를 바꾼 계산이다. 에어로졸과 기체 흡수 없이도 해면 반사 때문에 RAA에 따른 모양이 크게 달라진다. 아래 U 곡선은 RAA를 0--180도로 접을 때 어떤 정보가 사라지는지 보여준다. 이 곡선은 도식용 가상값이 아니라 함께 배포한 CSV를 그린 것이다. 방향별 반사도 $\rho_I$는 좁은 글린트 봉우리에서 1을 넘을 수 있다. 에너지 수지는 반구 적분 플럭스로 판단한다.
% Curves use actual OCRT C output, not illustrative fitted functions.
\begin{figure}[H]
\centering
\begin{tikzpicture}
\begin{axis}[
 width=15.3cm,height=6.5cm,
 xmin=0,xmax=360,ymin=0,
 xtick={0,90,180,270,360},
 xlabel={\FigLang{OCRT RAA (도)}{OCRT RAA (degrees)}},ylabel={$\rho_I$},
 grid=major,grid style={ocrtLine!65},axis line style={ocrtGray},
 ticklabel style={font=\sffamily\footnotesize},label style={font=\sffamily\small},
 legend style={font=\sffamily\footnotesize,at={(0.02,0.98)},anchor=north west,
   draw=none,fill=white,fill opacity=0.93,text opacity=1,cells={anchor=west}},
 clip=true]
\addplot[very thick,orange!85!black] table[x=raa_deg,y=rho_I_w1,col sep=comma]{figures/data/rayleigh_azimuth.csv};
\addlegendentry{\FigLang{직달 글린트 포함: 바람 1 m/s}{Direct glint included: wind 1 m/s}}
\addplot[very thick,ocrtTeal] table[x=raa_deg,y=rho_I_w5,col sep=comma]{figures/data/rayleigh_azimuth.csv};
\addlegendentry{\FigLang{직달 글린트 포함: 바람 5 m/s}{Direct glint included: wind 5 m/s}}
\addplot[thick,ocrtNavy,dashed] table[x=raa_deg,y=rho_I_decoupled,col sep=comma]{figures/data/rayleigh_azimuth.csv};
\addlegendentry{\FigLang{직달 글린트 제외: 바람 5 m/s}{Direct glint excluded: wind 5 m/s}}
\end{axis}
\end{tikzpicture}

\vspace{2mm}
\begin{tikzpicture}
\begin{axis}[
 width=15.3cm,height=5.2cm,
 xmin=0,xmax=360,xtick={0,90,180,270,360},
 xlabel={\FigLang{OCRT RAA (도)}{OCRT RAA (degrees)}},ylabel={$\rho_U$},
 grid=major,grid style={ocrtLine!65},axis line style={ocrtGray},
 ticklabel style={font=\sffamily\footnotesize},label style={font=\sffamily\small},
 scaled y ticks=false,yticklabel style={/pgf/number format/fixed,/pgf/number format/precision=3}]
\addplot[thick,ocrtGray,dotted,forget plot] coordinates {(0,0) (360,0)};
\addplot[very thick,ocrtTeal] table[x=raa_deg,y=rho_U_decoupled,col sep=comma]{figures/data/rayleigh_azimuth.csv};
\end{axis}
\end{tikzpicture}
\caption{\FigLang{실제 OCRT 계산으로 본 글린트 방향과 U의 거울 대칭. 555 nm, SZA=VZA=40도, 1013.25 hPa, \texttt{black\_fresnel\_ocean}, 에어로졸·6기체 흡수 없음, IPSS 끔이다. 위: 직달 글린트가 포함된 두 곡선의 최대값은 RAA=180도에 있다. 바람은 표면 경사 분포와 곡선 폭을 바꾼다. \texttt{--decouple-sunglint}는 직달 단일반사 항을 제외하며 해면의 모든 기여를 없애지는 않는다. 아래: 방위를 $\phi$에서 $360^\circ-\phi$로 바꾸면 U의 부호가 바뀐다. 360도 점은 그림을 닫기 위해 0도 값을 복사한 점이며 원래 LUT에는 없다.}{Glint direction and U mirror symmetry in actual OCRT calculations. Conditions: 555 nm, SZA=VZA=40 degrees, 1013.25 hPa, \texttt{black\_fresnel\_ocean}, no aerosol or six-gas absorption, IPSS off. Top: both curves including direct glint peak at RAA=180 degrees. Wind changes the surface-slope distribution and angular width. \texttt{--decouple-sunglint} removes the direct single-reflection term, while other surface contributions remain. Bottom: mirroring azimuth from $\phi$ to $360^\circ-\phi$ changes the sign of U. The plotted 360-degree endpoint repeats the 0-degree value for closure; it is absent from the native LUT.}}
\label{fig:measured-azimuth}
\end{figure}

\FloatBarrier

그림 자료는 저장소 최상위에서 다음 명령으로 재생성할 수 있다. 표준 라이브러리만 사용하는 Python 3가 필요하다. 원시 격자는 사례당 144행이고, 그림에는 VZA=40도인 72방위만 사용한다. 계산 조건·실행 파일 해시·대칭 검사는 메타데이터 JSON에 저장한다. 이 실행 검사는 생산 LUT의 수렴 정확도 전체를 검증하는 것은 아니다.
\par\noindent\begin{minipage}{\linewidth}
\begin{ManualCode}
python3 docs/manual/examples/generate_figure_data.py \
  --ocrt-root code/OCRT_C --exe code/OCRT_C/build/ocrt \
  --outdir task/manual_figure_data
\end{ManualCode}
\end{minipage}\par

\subsection{한 파장·SZA의 각도 격자}
\label{sec:auto:04_examples:3.2}

\par\noindent\begin{minipage}{\linewidth}
\begin{ManualCode}
./build/ocrt --surface black_fresnel_ocean --wind-speed 5 \
  --sza 40 --wavelength 555 --pressure 1013.25 \
  --decouple-sunglint "${GAS_OFF[@]}" \
  --lut-vza-step 30 --lut-vza-max 60 --lut-raa-step 90 \
  --output-full-grid ../../task/manual_examples/ray_w555_s40_p1013_w5.csv \
  2> ../../task/manual_examples/ray_w555_s40_p1013_w5.log
\end{ManualCode}
\end{minipage}\par

VZA는 \texttt{0,\allowbreak{}30,\allowbreak{}60\textdegree{}}, RAA는 \texttt{0,\allowbreak{}90,\allowbreak{}180,\allowbreak{}270\textdegree{}}로 12행이다. 360\textdegree{}는 0\textdegree{}와 중복되므로 포함되지 않는다. RAA 0--180\textdegree{}로 단순 접으면 U의 부호 정보를 잃을 수 있다. \texttt{-\allowbreak{}-\allowbreak{}lut-\allowbreak{}raa-\allowbreak{}step}은 360을 정확하게 나누는 값으로 선택한다.

기본 2.5\textdegree{}/VZA 최대85\textdegree{}라면 한 사례는 \texttt{(floor(85/\allowbreak{}2.\allowbreak{}5)+1) \ensuremath{\times} floor(36\allowbreak{}0/\allowbreak{}2.\allowbreak{}5) =\allowbreak{} 35\ensuremath{\times}144 =\allowbreak{} 5,\allowbreak{}040행}이다. 파장·SZA·기압·풍속 조건 수를 여기에 곱해 총 저장 규모를 계산한다. 격자는 출력 방향의 표본이고 \texttt{-\allowbreak{}-\allowbreak{}n-\allowbreak{}mu} 등 내부 적분 해상도와 별개다.

\subsection{여러 파장·SZA·기압·풍속}
\label{sec:auto:04_examples:3.3}
아래는 \textbf{저장소 최상위에서} 실행한다. 먼저 \texttt{-\allowbreak{}-\allowbreak{}run} 없이 16개 사례의 계획을 생성한다. 같은 명령에 \texttt{-\allowbreak{}-\allowbreak{}run}을 추가하면 실행한다.


\par\noindent\begin{minipage}{\linewidth}
\begin{ManualCode}
python3 docs/manual/examples/make_c_luts.py rayleigh \
  --ocrt-root code/OCRT_C --exe code/OCRT_C/build/ocrt \
  --outdir task/rayleigh_lut_demo \
  --wavelengths 443,555 --szas 20,60 \
  --pressures 900,1013.25 --winds 3,7 \
  --vza-step 30 --vza-max 60 --raa-step 90
\end{ManualCode}
\end{minipage}\par

이 예시는 16\ensuremath{\times}12=192행이며 생산용 권장 해상도가 아니다. \texttt{manifest.\allowbreak{}json}에는 실행 인자, 각 축, 기체·선글린트·IPSS 설정, 실행 파일 해시가 남는다. \texttt{jobs.\allowbreak{}csv}로 사례 파일과 입력 조건을 연결한다.

실제 위성 파장은 센서 밴드 정의를 기준으로 지정한다. 이 예시의 443/555 nm가 특정 센서의 전체 밴드 구성을 뜻하지 않는다. OCRT의 \texttt{-\allowbreak{}-\allowbreak{}wavelength}는 단색 계산이다. 넓은 밴드의 효과가 필요하면 센서 분광응답·태양 스펙트럼·흡수 구조를 고려한 별도 분광 적분 절차를 검증한다. 중심파장 한 번의 계산을 밴드 적분 결과로 표시하지 않는다.

\subsection{기압 보간·각도 보간}
\label{sec:auto:04_examples:3.4}
Rayleigh 광학두께는 기압에 비례하지만 다중산란과 해면을 포함한 최종 반사도가 기압에 정확히 선형인 것은 아니다. 결과 전체에 단순 \texttt{P/\allowbreak{}1013.\allowbreak{}25}를 곱해 기압 축을 대체하지 않는다. SZA/VZA/RAA 보간도 독립 계산점으로 오차를 평가한다. 큰 천정각이나 선글린트 부근은 더 조밀한 표본이 필요할 수 있다.

검증할 항목은 다음과 같다.

\begin{enumerate}
\item 파일별 행 수와 실제 헤더가 예상대로인지 확인한다.
\item \texttt{rho\_\allowbreak{}I/\allowbreak{}Q/\allowbreak{}U}에 NaN/Inf가 없는지 확인한다. Q/U의 음수는 정상적인 편광 정보일 수 있다.
\item 동일 모드에서 단일 방향 계산과 격자값을 대조한다. 두 경로의 내부 절점·기본값을 같게 맞추고 천저는 별도로 확인한다.
\item 기본 수치 설정과 더 세밀한 설정을 비교해 목적에 맞는 오차 한도를 정한다.
\item 기체와 직달 선글린트의 포함 여부, RAA와 편광 기준면을 대기보정 코드와 맞춘다.
\end{enumerate}

\Needspace{17\baselineskip}
\section{IPSS 구면 보정이 있는 Rayleigh LUT}
\label{sec:auto:04_examples:4}
\ref{sec:auto:04_examples:3}절의 명령에 \texttt{-\allowbreak{}-\allowbreak{}pssa}를 추가한다. 예를 들어:


\par\noindent\begin{minipage}{\linewidth}
\begin{ManualCode}
./build/ocrt --surface black_fresnel_ocean --wind-speed 5 \
  --sza 75 --wavelength 555 --pressure 1013.25 \
  --decouple-sunglint --pssa "${GAS_OFF[@]}" \
  --lut-vza-step 10 --lut-vza-max 80 --lut-raa-step 30 \
  --output-full-grid ../../task/manual_examples/ray_ipss_s75.csv \
  2> ../../task/manual_examples/ray_ipss_s75.log
\end{ManualCode}
\end{minipage}\par

옵션 이름은 과거 이름을 유지하지만 실제 알고리즘은 IPSS 하나다. SZA/VZA는 지표 화소 기준이다. 평면평행 LUT와 IPSS LUT는 메타데이터와 파일 이름으로 구분한다. \texttt{-\allowbreak{}-\allowbreak{}pssa-\allowbreak{}mode}는 삭제됐으며 \texttt{-\allowbreak{}-\allowbreak{}surface ocean -\allowbreak{}-\allowbreak{}water-\allowbreak{}model .\allowbreak{}.\allowbreak{}.\allowbreak{} -\allowbreak{}-\allowbreak{}pssa}는 지원되지 않는다. 대기·흑해양의 IPSS 지원을 결합 해양으로 확대해서 해석하지 않는다.

\Needspace{17\baselineskip}
\section{에어로졸 + Rayleigh + 해면}
\label{sec:auto:04_examples:5}
\texttt{inputs/\allowbreak{}M80C.\allowbreak{}mie}를 설치한 뒤:


\par\noindent\begin{minipage}{\linewidth}
\begin{ManualCode}
./build/ocrt --surface black_fresnel_ocean --wind-speed 5 \
  --sza 40 --vza 30 --raa 90 --wavelength 443 --pressure 1013.25 \
  --mie inputs/M80C.mie --aod-555 0.1 \
  --decouple-sunglint "${GAS_OFF[@]}" \
  > ../../task/manual_examples/aer443.txt \
  2> ../../task/manual_examples/aer443.log
\end{ManualCode}
\end{minipage}\par

AOD는 555 nm에서 0.1이며, 443 nm의 AOD는 Mie 소광비로 환산한다. \texttt{AOD\_\allowbreak{}ref\_\allowbreak{}nm}과 \texttt{AOD\_\allowbreak{}band}를 확인한다. 865 nm를 기준으로 하려면 \texttt{-\allowbreak{}-\allowbreak{}aod-\allowbreak{}555 0.\allowbreak{}1}을 \texttt{-\allowbreak{}-\allowbreak{}aod-\allowbreak{}865 0.\allowbreak{}1}로 바꾼다. 두 옵션을 함께 지정하지 않는다. AOD=0인 사례에서는 \texttt{-\allowbreak{}-\allowbreak{}mie} 등 에어로졸 설정도 뺀다.

에어로졸이 없는 결과를 뺀 차이에는 에어로졸과 분자·해수면의 상호작용도 포함된다. 따라서 이 차이가 단일산란 에어로졸 항과 같다고 가정하지 않는다. 모델명의 물리적 의미는 자료 배포처의 정의를 확인한다. Mie 모델을 바꿀 때는 모델명과 파일 해시도 기록한다.

\Needspace{17\baselineskip}
\section{기체 흡수와 대기 프로파일}
\label{sec:auto:04_examples:6}
일반 단일 기하와 명시적 전체 격자 경로에서는 기체 흡수가 기본으로 켜져 있다. 흡수까지 포함하는 다음 예시에는 \texttt{GAS\_\allowbreak{}OFF}를 붙이지 않는다.


\par\noindent\begin{minipage}{\linewidth}
\begin{ManualCode}
./build/ocrt --surface black_fresnel_ocean --wind-speed 5 \
  --sza 40 --vza 30 --raa 90 --wavelength 760 --pressure 1013.25 \
  --atm-profile tropical --gas-column-h2o 3.0 --gas-column-o3 300 \
  > ../../task/manual_examples/gas760.txt \
  2> ../../task/manual_examples/gas760.log
\end{ManualCode}
\end{minipage}\par

H2O의 단위는 g/cm\textsuperscript{2}, O3/NO2는 DU, O2/CO2/CH4는 \textbf{molecules/cm\textsuperscript{2}}다. 기존 도움말의 \texttt{mol/\allowbreak{}cm\textsuperscript{2}} 표시는 분자 수를 뜻한다. 이를 화학 단위인 mol로 해석하면 아보가드로수만큼의 배율 오류가 생긴다. 프로파일의 형태를 유지하면서 기체 총량을 조정하며, 세부 조건은 \hyperref[chap:03_options]{\ref*{chap:03_options}장}을 참고한다.

기체 흡수를 포함한 Rayleigh 배경을 만들 때는 \ref{sec:auto:04_examples:3}절의 기체 흡수를 끈 자료와 별도 데이터 집합으로 관리한다. 기존 \texttt{-\allowbreak{}-\allowbreak{}batch}는 일반 기체 초기화보다 먼저 분기하므로 이 용도에 사용하지 않는다.

\Needspace{17\baselineskip}
\section{결합 해양: 순수 해수와 CDOM}
\label{sec:auto:04_examples:7}
먼저 대기를 포함한 순수 해수 계산이다.


\par\noindent\begin{minipage}{\linewidth}
\begin{ManualCode}
./build/ocrt --surface ocean --water-model ocrt \
  --ocrt-chl 0 --ocrt-tsm 0 --ocrt-adom440 0 --wind-speed 3 \
  --sza 40 --vza 30 --raa 90 --wavelength 555 --pressure 1013.25 \
  > ../../task/manual_examples/pure_water.txt \
  2> ../../task/manual_examples/pure_water.log
\end{ManualCode}
\end{minipage}\par

CDOM의 기여를 추가하려면 \texttt{-\allowbreak{}-\allowbreak{}ocrt-\allowbreak{}adom440 0.\allowbreak{}1 -\allowbreak{}-\allowbreak{}ocrt-\allowbreak{}adom-\allowbreak{}slope 0.\allowbreak{}014}를 사용한다. 다음은 같은 조건의 전체 격자 예시다. CSV에 TOA 반사도, 수면 위 \texttt{Rrs}와 수면 아래 \texttt{rrs}, IOP, 투과 진단량 등을 함께 저장한다.


\par\noindent\begin{minipage}{\linewidth}
\begin{ManualCode}
./build/ocrt --surface ocean --water-model ocrt \
  --ocrt-chl 0 --ocrt-tsm 0 --ocrt-adom440 0.1 \
  --ocrt-adom-slope 0.014 --wind-speed 3 \
  --sza 40 --wavelength 555 --pressure 1013.25 \
  --lut-vza-step 30 --lut-vza-max 60 --lut-raa-step 90 \
  --output-full-grid ../../task/manual_examples/cdom555.csv \
  2> ../../task/manual_examples/cdom555.log
\end{ManualCode}
\end{minipage}\par

대기가 없는 수중 경계 검증에는 기압 0, 에어로졸 없음, 6기체 총량 0을 함께 지정한다. 이 경우에도 태양 조명과 공기--해수 경계는 남아 있다. 실제 해역의 관측값과 비교하는 조건과는 구분한다.

\Needspace{17\baselineskip}
\section{TSM·클로로필·쇄설물}
\label{sec:auto:04_examples:8}
필요한 Mie/IOP 자료를 설치한 뒤 실행하는 구성성분 모델 예시다.


\par\noindent\begin{minipage}{\linewidth}
\begin{ManualCode}
./build/ocrt --surface ocean --water-model ocrt \
  --ocrt-chl 1 --ocrt-tsm 0.5 --ocrt-adom440 0.02 \
  --ocrt-tsm-species red_clay --ocrt-detritus-a440 0.01 \
  --ocrt-adom-slope 0.014 --ocrt-detritus-slope 0.0109 \
  --water-temperature 20 --water-salinity 35 --wind-speed 3 \
  --sza 40 --vza 30 --raa 90 --wavelength 555 --pressure 1013.25 \
  > ../../task/manual_examples/constituents555.txt \
  2> ../../task/manual_examples/constituents555.log
\end{ManualCode}
\end{minipage}\par

Chl의 단위는 mg/m\textsuperscript{3}, 무기 TSM은 g/m\textsuperscript{3}, aDOM과 쇄설물의 기준 흡수는 m\textsuperscript{-}\textsuperscript{1}다. \textbf{현재 OCRT 구성성분 경로에서는 EAP 식물플랑크톤 산란이 비활성화되어 있으므로 Chl을 지정해도 해당 산란이 켜지지 않는다.} 일반 예시에 \texttt{-\allowbreak{}-\allowbreak{}ocrt-\allowbreak{}phyto-\allowbreak{}group}을 추가하지 않는다. 출력에서 흡수·산란의 성분별 값을 확인하고, 이 제약이 시뮬레이션 목적에 미치는 영향을 판단한다.

\Needspace{17\baselineskip}
\section{CCRR 입력 변환 경로}
\label{sec:auto:04_examples:9}

\par\noindent\begin{minipage}{\linewidth}
\begin{ManualCode}
./build/ocrt --surface ocean --water-model ccrr \
  --ccrr-chl 1 --ccrr-tsm 0 --ccrr-adom440 0.02 \
  --wind-speed 3 --sza 40 --vza 30 --raa 90 \
  --wavelength 555 --pressure 1013.25 \
  > ../../task/manual_examples/ccrr555.txt \
  2> ../../task/manual_examples/ccrr555.log
\end{ManualCode}
\end{minipage}\par

CCRR 입력을 IOP 등으로 변환한 뒤 OCRT로 계산하는 경로다. 별도의 복사전달 솔버를 실행하는 옵션은 아니다. \texttt{-\allowbreak{}-\allowbreak{}ocrt-\allowbreak{}*}와 \texttt{-\allowbreak{}-\allowbreak{}ccrr-\allowbreak{}*}의 구성성분 옵션을 섞지 않는다. 같은 Chl 값을 넣어도 OCRT 구성성분 모델과 같은 IOP가 만들어진다고 보장할 수 없다.

\Needspace{17\baselineskip}
\section{전체 IOP 직접 지정}
\label{sec:auto:04_examples:10}

\par\noindent\begin{minipage}{\linewidth}
\begin{ManualCode}
./build/ocrt --surface ocean --water-model iop \
  --iop-a 0.1 --iop-b 0.2 --iop-bb 0.005 --wind-speed 3 \
  --sza 40 --vza 30 --raa 90 --wavelength 555 --pressure 1013.25 \
  > ../../task/manual_examples/iop555.txt \
  2> ../../task/manual_examples/iop555.log
\end{ManualCode}
\end{minipage}\par

세 값은 순수 해수 등을 포함한 \textbf{전체} a, b, bb이며 단위는 모두 m\textsuperscript{-}\textsuperscript{1}다. \texttt{a\ensuremath{\geq}0,\allowbreak{} b>0,\allowbreak{} 0<bb<0.\allowbreak{}5b}를 만족해야 한다. bb/b만으로 위상함수의 전체 형태가 정해지지는 않는다. 비교 목적에 맞게 위상함수·보간·절단 방법을 선택하고, 필요하면 \texttt{-\allowbreak{}-\allowbreak{}iop-\allowbreak{}phase-\allowbreak{}lut} 또는 \texttt{-\allowbreak{}-\allowbreak{}iop-\allowbreak{}mie-\allowbreak{}phase}를 사용한다. \hyperref[chap:03_options]{\ref*{chap:03_options}장}의 옵션 충돌 조건을 지킨다. 과학 참조 계산의 기본 Mie 전방 피크 절단은 OFF이며, 속도를 높이려고 임의로 절단 옵션을 추가하지 않는다.

\Needspace{17\baselineskip}
\section{시뮬레이션 자료의 체계적 생성}
\label{sec:auto:04_examples:11}
저장 단위는 하나의 대기·해양 상태에 대한 각도 격자다. 관측량을 담은 행과 상태를 설명하는 표를 연결해 관리한다. 다음 예시는 2파장\ensuremath{\times}2SZA\ensuremath{\times}2CDOM=8상태, 상태당 12방향이다. \textbf{이 명령은 다시 저장소 최상위에서 실행하며, 출력 경로와 실행 파일 경로도 그 위치를 기준으로 지정한다.}


\par\noindent\begin{minipage}{\linewidth}
\begin{ManualCode}
python3 docs/manual/examples/make_c_luts.py simulation \
  --ocrt-root code/OCRT_C --exe code/OCRT_C/build/ocrt \
  --outdir task/simulation_demo \
  --wavelengths 443,555 --szas 20,60 --winds 3 \
  --adom440 0,0.1 --vza-step 30 --vza-max 60 --raa-step 90
\end{ManualCode}
\end{minipage}\par

계획을 확인한 뒤 \texttt{-\allowbreak{}-\allowbreak{}run}을 추가한다. 기체는 기본 usstd76 프로파일을 사용하고 에어로졸은 없으며 Chl/TSM은 0이다. 에어로졸 축을 추가하려면 \texttt{-\allowbreak{}-\allowbreak{}aods 0,\allowbreak{}0.\allowbreak{}1 -\allowbreak{}-\allowbreak{}mie inputs/\allowbreak{}M80C.\allowbreak{}mie}, 기체 흡수를 끄려면 \texttt{-\allowbreak{}-\allowbreak{}gas-\allowbreak{}off}를 사용한다. 이 보조 도구에는 Chl/TSM 축이 없다. 해당 성분을 바꾸려면 \ref{sec:auto:04_examples:8}절의 명시적 인자를 사용하는 별도 사례표를 만든다.

학습 자료 등으로 사용할 때도 입력 농도와 IOP, TOA 반사도와 수면 \texttt{Rrs}를 각각 별도 열로 유지한다. 실패하거나 수렴하지 않은 행을 0으로 채우지 않는다. \texttt{water\_\allowbreak{}converged}, 실행 종료코드, NaN/Inf, 투과량 유효성 플래그를 확인한다. 같은 대기·해양 상태의 인접 각도만 무작위로 학습·평가 집합에 나누면 자료의 독립성을 과대평가할 수 있다. 사용 목적에 맞춰 대기·해양 상태 단위의 평가 집합을 설계한다.

\Needspace{17\baselineskip}
\section{C 기본 해양 전체 격자 배치}
\label{sec:auto:04_examples:12}
\ref{sec:auto:04_examples:1}절과 같은 C 작업 폴더에서 시작한다. 먼저 다음 CSV를 \texttt{.\allowbreak{}.\allowbreak{}/\allowbreak{}.\allowbreak{}.\allowbreak{}/\allowbreak{}task/\allowbreak{}manual\_\allowbreak{}examples/\allowbreak{}ocean\_\allowbreak{}jobs.\allowbreak{}csv}에 UTF-8 BOM 없이 저장한다.


\par\noindent\begin{minipage}{\linewidth}
\begin{ManualCode}
out,wavelength,sza,wind,ocrt_adom440
../../task/manual_examples/b443.csv,443,40,3,0.1
../../task/manual_examples/b555.csv,555,40,3,0.1
\end{ManualCode}
\end{minipage}\par


\par\noindent\begin{minipage}{\linewidth}
\begin{ManualCode}
./build/ocrt --surface ocean --water-model ocrt \
  --ocrt-chl 0 --ocrt-tsm 0 --ocrt-adom440 0.1 --wind-speed 3 \
  --sza 40 --wavelength 443 --pressure 1013.25 --n-water 1.34 \
  --lut-vza-step 30 --lut-vza-max 60 --lut-raa-step 90 \
  --batch-full-grid ../../task/manual_examples/ocean_jobs.csv \
  > ../../task/manual_examples/ocean_batch.txt \
  2> ../../task/manual_examples/ocean_batch.log
\end{ManualCode}
\end{minipage}\par

CSV에 없는 값이나 빈칸은 기본 CLI 값을 사용한다. CSV에 SZA와 파장이 있더라도 CLI에 기본 \texttt{-\allowbreak{}-\allowbreak{}sza/\allowbreak{}-\allowbreak{}-\allowbreak{}wavelength}를 지정한다. \texttt{out}은 사례마다 서로 다른 값으로 지정한다. 입력기는 완전한 범용 CSV 파서가 아니므로 필드 안에 쉼표·따옴표·줄바꿈을 넣지 않는다. 지원하는 모든 입력 열은 \hyperref[chap:08_input_formats]{\ref*{chap:08_input_formats}장 입력 파일 포맷}을 참고한다.

이 예시는 앞의 단일 상태 계산과 굴절률을 맞추기 위해 \texttt{-\allowbreak{}-\allowbreak{}n-\allowbreak{}water 1.\allowbreak{}34}를 명시한다. 현재 배치 구현은 행에 파장이 있고 CLI에서 굴절률을 지정하지 않으면 Quan--Fry 파장·수온·염분 의존 굴절률을 선택한다. 따라서 옵션을 생략한 단일 계산의 고정 굴절률 1.34와 다를 수 있다. 파장 의존 굴절률을 사용할 때는 모든 비교 경로에서 같은 선택을 적용한다.

이 배치 명령은 \texttt{-\allowbreak{}-\allowbreak{}output-\allowbreak{}full-\allowbreak{}grid}로 요청한 단일 상태 격자의 수중 96절점 기본값을 자동으로 계승하지 않는다. 현재 소스는 \texttt{-\allowbreak{}-\allowbreak{}output-\allowbreak{}full-\allowbreak{}grid}가 명시적으로 격자 모드를 켠 경우에 수중 기본값을 96으로 올리며, 위 배치 명령은 그 조건을 거치지 않아 일반 수중 기본값 48을 사용할 수 있다. 두 경로를 대조하거나 같은 수치 설정의 자료를 만들려면 \texttt{OCRT\_\allowbreak{}ADVANCED=\allowbreak{}1} 환경에서 \texttt{-\allowbreak{}-\allowbreak{}n-\allowbreak{}mu-\allowbreak{}water 96}을 양쪽에 명시한다. 고농도 조건에서는 별도로 수렴을 확인한다.

\texttt{OMP\_\allowbreak{}NUM\_\allowbreak{}THREADS}를 늘리면 행 단위 병렬 실행에 사용한다. 한 사례의 메모리 사용량을 확인한 뒤 스레드 수를 늘린다. 기압처럼 CSV에서 제어하지 않는 조건은 별도 CLI 실행으로 나눈다. \textbf{일반 \texttt{-\allowbreak{}-\allowbreak{}batch}와 \texttt{-\allowbreak{}-\allowbreak{}batch-\allowbreak{}full-\allowbreak{}grid}는 구현 경로가 다르다.} 전체 격자 배치에서 지원하는 기능을 일반 배치도 지원한다고 가정하지 않는다.

\Needspace{24\baselineskip}
\section{Python으로 CSV 읽기와 점검}
\label{sec:auto:04_examples:13}
Python 표준 라이브러리만으로 읽을 수 있다. 다음은 C 작업 폴더에서 \ref{sec:auto:04_examples:7}절의 CSV를 점검하는 예시다.


\par\noindent\begin{minipage}{\linewidth}
\begin{ManualCode}
import csv
import math
from pathlib import Path

path = Path('../../task/manual_examples/cdom555.csv')
with path.open(encoding='utf-8', newline='') as f:
    rows = list(csv.DictReader(f))
assert len(rows) == 12
for row in rows:
    for name in ('TOA_rho_I', 'Rrs_I', 'rrs_I', 'a_total', 'b_total'):
        assert math.isfinite(float(row[name])), (name, row)
    assert int(row['water_converged']) == 1, row
    if int(row['T_up_rt_valid']) != 1:
        print('Treat upward RT transmittance as invalid:', row['vza_deg'], row['raa_deg'])
print(rows[0]['wavelength_nm'], rows[0]['Rrs_I'])
\end{ManualCode}
\end{minipage}\par

\texttt{Rrs\_\allowbreak{}I}의 단위는 sr\textsuperscript{-}\textsuperscript{1}이며 \texttt{TOA\_\allowbreak{}rho\_\allowbreak{}I}는 무차원이다. \texttt{Rrs\_\allowbreak{}I}에 \ensuremath{\pi}를 곱한 값은 수면의 하향 복사조도로 정규화한 수면 반사량으로, TOA 반사도와는 다르다. 저장할 때 단위, 0+/0- 위치, 기준면, 기체·선글린트 조건을 함께 기록한다.

\Needspace{17\baselineskip}
\section{첫 계산: Bash와 PowerShell}
\label{sec:auto:04_examples:14}

\label{sec:example_first_run}

\textbf{실행할 때의 현재 디렉토리를 \texttt{code/\allowbreak{}OCRT\_\allowbreak{}C}로 둔다.} 실행 파일의 위치만 지정해도 내부의 \texttt{inputs/\allowbreak{}} 상대 경로가 자동으로 바뀌지는 않는다.

\Needspace{13\baselineskip}
Bash:


\par\noindent\begin{minipage}{\linewidth}
\begin{ManualCode}
export OMP_NUM_THREADS=1
mkdir -p ../../task/manual_examples
./build/ocrt --surface black --sza 40 --vza 30 --raa 90 \
  --wavelength 555 --pressure 1013.25 \
  --gas-column-h2o 0 --gas-column-o3 0 --gas-column-no2 0 \
  --gas-column-o2 0 --gas-column-co2 0 --gas-column-ch4 0 \
  > ../../task/manual_examples/rayleigh_single.txt \
  2> ../../task/manual_examples/rayleigh_single.log
\end{ManualCode}
\end{minipage}\par

\Needspace{13\baselineskip}
PowerShell:


\par\noindent\begin{minipage}{\linewidth}
\begin{ManualCode}
$env:OMP_NUM_THREADS = '1'
New-Item -ItemType Directory -Force ..\..\task\manual_examples | Out-Null
$gasOff = @('--gas-column-h2o','0','--gas-column-o3','0','--gas-column-no2','0',
            '--gas-column-o2','0','--gas-column-co2','0','--gas-column-ch4','0')
& .\build\ocrt_x86-64-v3.exe --surface black --sza 40 --vza 30 --raa 90 `
  --wavelength 555 --pressure 1013.25 @gasOff `
  > ..\..\task\manual_examples\rayleigh_single.txt `
  2> ..\..\task\manual_examples\rayleigh_single.log
if ($LASTEXITCODE -ne 0) { throw 'OCRT failed; inspect the log' }
\end{ManualCode}
\end{minipage}\par

이것은 검은 경계 위의 분자 산란 계산이며 해면 반사는 포함하지 않는다. 해양 대기보정용 해면 포함 Rayleigh LUT는 \hyperref[chap:04_examples]{부록~\ref*{chap:04_examples}}를 사용한다. 표준출력에는 처음 I/Q/U 세 숫자와 \texttt{key=\allowbreak{}value} 형식의 결과, 표준에러에는 진단이 나온다. \hyperref[chap:05_output_formats]{\ref*{chap:05_output_formats}장}에서 의미를 확인한다.

\Needspace{24\baselineskip}
\section{CSV 읽기: Rrs와 rrs 구분}
\label{sec:auto:04_examples:15}

\label{sec:example_csv_reader}

해양 격자 출력 파일을 읽어 수면 위와 수면 아래 반사도를 각각 가져오는 예시다. 아래 파일명을 실제 결과 파일 경로로 바꾼다.

\textbf{열 이름의 대소문자를 보존한다.} 해양 CSV의 \texttt{Rrs\_\allowbreak{}I}와 \texttt{rrs\_\allowbreak{}I}는 서로 다른 물리량이다. 대소문자를 무시하는 PowerShell \texttt{Import-\allowbreak{}Csv}는 이 헤더를 중복으로 판단해 읽기에 실패할 수 있다. Python의 \texttt{csv.\allowbreak{}DictReader} 또는 \texttt{pandas.\allowbreak{}read\_\allowbreak{}csv}처럼 대소문자를 구별하는 방법으로 읽고, 열 이름을 일괄 소문자로 바꾸지 않는다.

\par\noindent\begin{minipage}{\linewidth}
\begin{ManualCode}
import csv

with open("ocean_grid.csv", newline="", encoding="utf-8-sig") as f:
    for row in csv.DictReader(f):
        above_water = float(row["Rrs_I"])
        below_water = float(row["rrs_I"])
        # above_water: 0+; below_water: 0-; 둘 다 sr^-1
\end{ManualCode}
\end{minipage}\par

\Needspace{17\baselineskip}
\section{두 상태 해양 배치 실행}
\label{sec:auto:04_examples:16}

\label{sec:example_input_batch}

다음 CSV는 폴더를 미리 만든 뒤 \texttt{jobs\_\allowbreak{}adom.\allowbreak{}csv}로 저장한다. 파일 경로는 \texttt{code/\allowbreak{}OCRT\_\allowbreak{}C}를 기준으로 한다.


\par\noindent\begin{minipage}{\linewidth}
\begin{ManualCode}
out,wavelength,sza,wind,ocrt_adom440
../../task/manual_examples/adom_443_s30.csv,443,30,3,0.05
../../task/manual_examples/adom_555_s40.csv,555,40,5,0.10
\end{ManualCode}
\end{minipage}\par


\par\noindent\begin{minipage}{\linewidth}
\begin{ManualCode}
export OMP_NUM_THREADS=1
export OCRT_ADVANCED=1
./build/ocrt --surface ocean --water-model ocrt \
  --ocrt-chl 0 --ocrt-tsm 0 --ocrt-adom440 0.1 \
  --sza 30 --wavelength 555 --wind-speed 3 --n-water 1.34 \
  --n-mu-water 96 --lut-vza-step 30 --lut-vza-max 60 --lut-raa-step 90 \
  --batch-full-grid jobs_adom.csv
\end{ManualCode}
\end{minipage}\par

위 명령의 기체흡수는 기본 프로파일로 활성이다. 기체를 제외하려면 \hyperref[chap:03_options]{\ref*{chap:03_options}장의 여섯 컬럼 0 옵션}을 CLI에 모두 추가한다. 각 행의 예상 격자 크기는 VZA 3개\ensuremath{\times}RAA 4개=12행이다. 이 짧은 예시는 입력 구조를 설명하며, 실제 자료 생성에는 \hyperref[chap:04_examples]{부록~\ref*{chap:04_examples}}의 확인된 실행 절차를 함께 따른다.

\Needspace{17\baselineskip}
\section{IPSS 검증 스크립트 실행}
\label{sec:auto:04_examples:17}
\label{sec:example_ipss_gates}

Linux/Bash에서 \texttt{code/\allowbreak{}OCRT\_\allowbreak{}C}를 작업 폴더로 두고 다음을 실행한다. GCC와 Python 3가 필요하다. 기하·Rayleigh 수치 검사는 대형 Mie 자료 없이 실행할 수 있다.


\par\noindent\begin{minipage}{\linewidth}
\begin{ManualCode}
bash scripts/run_ipss_gates.sh
\end{ManualCode}
\end{minipage}\par

기대 종료 표시는 \texttt{IPSS GATES:\allowbreak{} ALL PASS}와 \texttt{IPSS REGRESSION:\allowbreak{} ALL PASS}다. 스크립트는 \texttt{build/\allowbreak{}ocrt}를 사용하므로 검사 대상 최신 파일을 그 이름으로 빌드해 둔다. Windows \texttt{.\allowbreak{}exe} 이름으로 자동 전환하지 않는다. 결합 해양+IPSS 거부 검사는 실패 여부만으로 판정하므로 자료 누락으로 실패한 경우와 구분하려면 실제 에러 문구까지 확인한다.
